\documentclass[10pt,a4paper]{amsart}
\usepackage[margin=19mm]{geometry}
\usepackage{amsmath,amssymb,mathtools}
\usepackage[scr=rsfs,cal=euler]{mathalfa}
\usepackage{xfrac}
\usepackage[unicode,hidelinks,bookmarks=false]{hyperref}
\newcommand{\ie}{i.e.\ }
\newcommand{\cf}{cf.\ }
\newcommand{\resp}{respectively\ }
\newcommand{\Subsection}{\subsection}
\providecommand{\nobreakdash}{\nobreak\hbox{-}}
\setlength{\emergencystretch}{2em}
\multlinegap0pt
\usepackage{bm}
\usepackage{amssymb}
\usepackage[all]{xy}
\usepackage{enumerate}
\usepackage{float}
\usepackage{tikz}
\newtheorem{thm}{Theorem}
\newtheorem{cor}{Corollary}
\newtheorem{lem}{Lemma}
\newcommand{\lt}{{}^\mathrm t\hspace{-0.5mm}}
\newcommand{\pmat}[1]{\begin{pmatrix} #1 \end{pmatrix}}
\title{Combinatorial structure of low degree rational curves on a smooth\\Hermitian surface}
\author{Norifumi Ojiro}
\date{Source: arXiv:2602.10842v1 (CC BY 4.0)}
\begin{document}
\maketitle

\noindent\emph{Source and licence.} Combinatorial structure of low degree rational curves on a smooth\\Hermitian surface. Version arXiv:2602.10842v1 (CC BY 4.0), distributed by arXiv under the Creative Commons Attribution 4.0 International licence, http://creativecommons.org/licenses/by/4.0/, as recorded in the arXiv licence metadata for this exact version and shown on the abstract page https://arxiv.org/abs/2602.10842v1. Full licence notice: \texttt{licenses/arxiv-2602.10842-CC-BY-4.0.txt}.\par\smallskip

\noindent\textit{Editorial scope.} Counted unit: the article's lem lemcl (source0.tex lines 256--262). The article's complete original proof of it is delivered with it (source0.tex lines 263--326, one \texttt{proof} environment of 64 lines). No mathematics is written, repaired or re-derived: the statement and the proof are the article's own lines, in the article's own notation, and no retained line is substituted or re-flowed.  Retained uncounted as the source-local context the proof needs: lines 231--233; lines 250--252.  Nothing was omitted from any retained span, so the package carries no omission record.  No retained line contains a citation, so the wrapper carries no bibliography and no bibliography entry.  No retained line contains a figure environment, an image-inclusion command or drawing code, so no image is delivered and the package declares no asset dependency.  Editorial formatting only, in addition: an AMS \texttt{amsart} preamble that restates the article's own declarations and macros used by the retained lines, namely the environments \texttt{thm}, \texttt{cor}, \texttt{lem} on separate counters, together with the standard packages those lines need.  The retained environments print in this document as Theorem 1, Corollary 1, Lemma 1 in order of appearance, whereas the article prints the same items with the numbering of its own full document.  The complete native source of the article as distributed in the arXiv e-print for this exact version is bundled unchanged as \texttt{sources/194433/source0.tex}.
\begin{thm}\label{thmcl1}
If $O={\rm Aut}(X_I)C_{F_J}$ then for each $\bm{p}\in X_I(\mathbb{F}_{q^2})$ there are $C\in O$ and $L\in\mathcal{L}(X_I)$ such that $\bm{p}\in C\cap L$.
\end{thm}

\begin{cor}\label{corcl}
Let $O$ be as in Theorem \ref{thmcl1}. Then there is $L\in\mathcal{L}(X_I)$ such that $C_{F_J}\cap L\neq\emptyset$.
\end{cor}

\begin{lem}\label{lemcl}
%Let $C\in O$ and $L\in\mathcal{L}(X_I)$. 
If $O={\rm Aut}(X_I)C_{F_J}$ then
%\begin{center}
$|C(\mathbb{F}_{q^2})\cap L(\mathbb{F}_{q^2})|\le1$ for any $C\in O$ and $L\in\mathcal{L}(X_I)$.
%\end{center}
\end{lem}

\begin{proof}
%If $C\cap L=\emptyset$ then $|C(\mathbb{F}_{q^2})\cap L(\mathbb{F}_{q^2})|=0$. Further, by the assumption and 
By the transitivity on $\mathcal{R}(X_I)$ and Corollary \ref{corcl}, it suffices to prove that $|C_{F_J}(\mathbb{F}_{q^2})\cap L(\mathbb{F}_{q^2})|=1$ for $L\in\mathcal{L}(X_I)$ such that $C_{F_J}\cap L\neq\emptyset$.
%one may replace $C$ by $C_{F_J}$. 
%Put $L=L_G$ for some $G\in\mathcal{M}_{4,2}$. 
%Let us show contradiction by 
Assume that $|C_{F_J}(\mathbb{F}_{q^2})\cap L(\mathbb{F}_{q^2})|\ge2$. Put $L=L_G$, where $G\in\mathcal{M}_{4,2}(\mathbb{F}_{q^2})$ is of rank $2$. Then, since $\mathbb{P}^1(k)\rightarrow C_{F_J}$ is injective, there are $\lt(a_1,b_1),\lt(a_2,b_2)\in\mathbb{P}^1(k)$ with  $\lt(a_1,b_1)\neq\lt(a_2,b_2)$ such that 
% $(c_1,d_1)\neq(c_2,d_2)$ and
\begin{equation}\label{lemcl1}
F_J\pmat{{a_1}^{q+1}&{a_2}^{q+1}\\{a_1}^qb_1&{a_2}^qb_2\\a_1{b_1}^q&a_2{b_2}^q\\{b_1}^{q+1}&{b_2}^{q+1}}=G.
%\ {\rm for}\ \pmat{a_1\\b_1}\neq\pmat{a_2\\b_2},
%\pmat{c_1&c_2\\d_1&d_2}.
\end{equation}
%and so
Since $L_G\subset X_I$ we have
$$
\pmat{{a_1}^{q+1} & {a_1}^q b_1 & a_1{ b_1}^q & {b_1}^{q+1} \\ {a_2}^{q+1} & {a_2}^q b_2 & a_2 {b_2}^q & {b_2}^{q+1}}\lt F_J{F_J}^{(q)}\pmat{{a_1}^{q+1}&{a_2}^{q+1}\\{a_1}^qb_1&{a_2}^qb_2\\a_1{b_1}^q&a_2{b_2}^q\\{b_1}^{q+1}&{b_2}^{q+1}}^{(q)}=\pmat{0&0\\0&0},
$$
or equivalently,
$$
\begin{cases}
-({a_1}^qb_2)^{q+1}+{a_1}^{q^2}{b_1}^q{a_2}^qb_2+{a_1}^q{b_1}^{q^2}a_2{b_2}^q-({b_1}^qa_2)^{q+1}=0,\\
-(b_1{a_2}^q)^{q+1}+{a_1}^{q}{b_1}{a_2}^{q^2}{b_2}^q+{a_1}{b_1}^{q}{a_2}^q{b_2}^{q^2}-(a_1{b_2}^q)^{q+1}=0.
\end{cases}
$$
From these equations, we have 
\begin{center}
$a_1=0\iff a_2=0$ and $b_1=0\iff b_2=0$, 
\end{center}
and so $a_1b_1a_2b_2\neq0$. Therefore without loss of generality, one may put $b_1=b_2=1$ and $a_2=ca_1\neq0$ for $c\neq0,1$. Then $a_1,a_2\in{\mathbb{F}_{q^2}}^\times$ by \eqref{lemcl1}, and furthermore the equations are simplified to
\begin{equation*}%\label{lemcl2}
%\pmat{c^{q^2}-c^{q^2+q}&c^q-1\\c^q-1&c-c^{q+1}}\pmat{{a_1}^{q^2+q}\\{a_1}^{q+1}}=\pmat{0\\0}.
\begin{cases}
(c^{q^2}-c^{q^2+q}){a_1}^{q^2+q}+(c^q-1){a_1}^{q+1}=0,\\
(c^q-1){a_1}^{q^2+q}+(c-c^{q+1}){a_1}^{q+1}=0.
\end{cases}
\end{equation*}
%It is easy to see that the rank of the coefficient matrix $M$ of the above equation is $1\le2$, and so the determinant is zero. Hence we have
%$$
%%c^{q^2+1}-2c^{q^2+q+1}+c^{q^2+2q+1}-c^{2q}+2c^q-1
%{\rm det}(M)=(c^q-1)^2(c^{q^2+1}-1)=0,
%$$
%and so $c^{q^2+1}=1$. 
%Put $a={a_1}^{q^2+q}/{a_1}^{q+1}={a_1}^{q^2-1}$. 
Hence
% by \eqref{lemcl2}, 
we have
%$$
%(c^q-1){a_1}^{q^2-1}+(c-c^{q+1})=0,
%$$
%or equivalently,
\begin{center}
$(c^{q^2}{a_1}^{q^2-1}-1)(c^q-1)=0$ and $(c-{a_1}^{q^2-1})(c^{q}-1)=0$,
\end{center}
%where $a={a_1}^{q^2+q}/{a_1}^{q+1}$. 
and thus $c={a_1}^{q^2-1}$ and $c^{q^2}{a_1}^{q^2-1}=1$. However ${a_1}^{q^2-1}=1$ since $a_1\in{\mathbb{F}_{q^2}}^\times$, and so $c=1$
%, and consequently $a=c$ by $c^{q^2+1}=1$. 
%Since $a\in{\mathbb{F}_{q^2}}^{\times}$ by \eqref{lemcl1}, one has $c^{q^2-1}=a^{q^2-1}=1$. Combining with $c^{q^2+1}=1$ we have $c^2=1$, thus $c=-1$ and $q$ is odd, and then ${a_1}^{q^2-1}=a=c=-1$. Therefore the left hand side of \eqref{lemcl1} is of the form
%$$
%F_J\pmat{a_1^{q+1}&a_1^{q+1}\\{a_1}^q&-a_1^q\\a_1&-a_1\\1&1},
%$$
%and it belongs to $\mathcal{M}_{4,2}(\mathbb{F}_{q^2})$. Hence it must be true that ${a_1}^{q^2-1}=1$, 
a contradiction, thus the proof has completed.
\end{proof}

\end{document}
